\documentclass[16pt,a4paper]{article} % ใช้ขนาดฟอนต์ 16pt ตามมาตรฐานข้อสอบไทย
\usepackage{geometry}
\geometry{top=2.5cm, bottom=2.5cm, left=2.2cm, right=2.2cm}

% --- การตั้งค่าภาษาไทยและฟอนต์ ---
\usepackage{fontspec}
\setmainfont{TH Sarabun New} % สำหรับโปรแกรมในเครื่องคอมพิวเตอร์ที่มีฟอนต์นี้
% หากใช้ Overleaf แนะนำให้ดาวน์โหลดฟอนต์ TH Sarabun New มาใส่ในโปรเจกต์ หรือเปลี่ยนเป็น \setmainfont{Sarabun}

\usepackage{amsmath,amssymb}
\usepackage{graphicx}
\usepackage{fancyhdr}
\usepackage{lastpage}

% เพิ่มระยะห่างระหว่างบรรทัดเล็กน้อย เพื่อไม่ให้สระและวรรณยุกต์ภาษาไทยซ้อนทับกัน
\linespread{1.25}

% --- การตั้งค่าหัวกระดาษ (Header) แสดงเลขหน้า "หน้า X/Y" ---
\pagestyle{fancy}
\fancyhf{} % ล้างค่าเริ่มต้นทั้งหมด
\fancyhead[R]{\fontsize{12pt}{12pt}\selectfont หน้า \thepage} % จัดให้ชิดขวาตามข้อสอบจริง
\renewcommand{\headrulewidth}{0pt} % เอาเส้นคั่นแนวนอนออก
\renewcommand{\footrulewidth}{0pt}

% --- ตัวนับเลขข้อแบบรันต่อเนื่อง (Continuous Question Counter) ---
\newcounter{qnum}
\newcommand{\question}{\refstepcounter{qnum}\par\medskip\noindent\textbf{\theqnum. }}

% --- คำสั่งสำหรับสร้างหัวข้อแต่ละตอน ---
\newcommand{\examsection}[1]{%
  \vspace{0.6cm}
  \noindent\textbf{\large #1}
  \vspace{0.3cm}\par
}

% --- ระบบจัดรูปแบบช้อยส์ (ก. ข. ค. ง.) ---
% 1. ช้อยส์แบบแนวนอนเรียงแถวเดียว (เหมาะสำหรับตัวเลขสั้นๆ)
\newcommand{\choiceshorizontal}[4]{%
  \nopagebreak
  \noindent\makebox[0.22\textwidth][l]{ก. #1}%
  \makebox[0.22\textwidth][l]{ข. #2}%
  \makebox[0.22\textwidth][l]{ค. #3}%
  \makebox[0.22\textwidth][l]{ง. #4}\par\medskip
}

% 2. ช้อยส์แบบตาราง 2x2 (เรียง ก. ข. อยู่บน ค. ง. อยู่ล่าง เหมาะสำหรับข้อความสั้น-ปานกลาง)
\newcommand{\choicesgrid}[4]{%
  \nopagebreak
  \noindent\begin{tabular}{@{}p{0.45\textwidth}p{0.45\textwidth}}
    ก. #1 & ข. #2 \\
    ค. #3 & ง. #4
  \end{tabular}\par\medskip
}

% 3. ช้อยส์แบบเรียงดิ่งลงมาแถวละข้อ (เหมาะสำหรับช้อยส์ที่เป็นประโยคยาวๆ)
\newcommand{\choicesvertical}[4]{%
  \nopagebreak
  \noindent\hspace*{1em} ก. #1 \\
  \noindent\hspace*{1em} ข. #2 \\
  \noindent\hspace*{1em} ค. #3 \\
  \noindent\hspace*{1em} ง. #4 \par\medskip
}

\begin{document}

% ========= part 1 ============
\examsection{ตอนที่ 1 คณิตศาสตร์ จำนวน 30 ข้อ (ข้อ 1--30) ข้อละ 1 คะแนน}

\question 
\newline
\choiceshorizontal{749}{901}{911}{897}

% ตัวอย่างโจทย์คณิตศาสตร์ธรรมดา + ช้อยส์แนวนอน
\question จงหาผลบวกของจำนวนเต็มบวก $n$ ทั้งหมด ที่สอดคล้องกับ $1 < n < 1000$ โดยที่เศษจากการหารด้วย 7, 10 และ 13 คือ 1
\newline
\choiceshorizontal{749}{901}{911}{897}

% ตัวอย่างโจทย์คณิตศาสตร์ + ช้อยส์ตาราง 2x2
\question กำหนด $f(x)=2(ax)^{2}-ax-1$ และ $g(x)=ax-1$ โดยที่ $a>0$ ถ้า $\frac{f}{g}(2a)=2$ แล้ว จงหาค่าของ $f\circ g^{-1}(a)$ เมื่อ $x\ne\frac{1}{a}$
\newline
\choicesgrid{$-1$}{$\frac{1}{2}$}{$2$}{$3$}

\pagebreak
% ========= part 2 ============
\examsection{ตอนที่ 2 วิทยาการคำนวณ จำนวน 25 ข้อ (ข้อ 31--55) ข้อละ 1 คะแนน}

% ตัวอย่างโจทย์แบบไม่มีรูปภาพ (ข้อสอบเขียนโปรแกรม/Code Snippet)
\question โปรแกรมนี้พิมพ์คำตอบอะไร
\begin{verbatim}
x = 1
while x < 5:
    x = x + 3
    if x % 3 == 0:
        continue
    x = x - 4
print(x)
\end{verbatim}
\choiceshorizontal{4}{5}{6}{7}

% ตัวอย่างโจทย์แบบมีรูปภาพอยู่ตรงกลาง (ใช้ \framebox แทนรูปจำลองเพื่อให้โค้ดคอมไพล์ผ่านทันที)
\question จากรูปวงกลมที่มีจุด $O$ เป็นจุดศูนย์กลางและมีรัศมีเท่ากับ 2 หน่วยด้านล่างนี้ จงหาพื้นที่ของรูปสามเหลี่ยม $\triangle ABC$
\begin{center}
    % เมื่อใช้งานจริงให้เปลี่ยนบรรทัดล่างนี้เป็นบรรทัดเรียกรูปภาพ เช่น: \includegraphics[width=0.35\textwidth]{image1.png}
    \framebox[5cm]{\vbox to 4cm{\vfill\hbox to 5cm{\hfill [บล็อกจำลองรูปภาพ] \hfill}\vfill}}
\end{center}
\choicesgrid{2}{$\sqrt{3}$}{$\sqrt{3}-1$}{$2\sqrt{3}-2$}

% ตัวอย่างโจทย์แบบมีรูปภาพอยู่ด้านขวา และข้อความอยู่ด้านซ้าย (ใช้ minipage)
\question พิจารณารูปสามเหลี่ยมมุมฉากที่กำหนดให้ต่อไปนี้ 
\nopagebreak
\begin{minipage}{0.55\textwidth}
    กำหนดรูปสามเหลี่ยม $ABC$ มีมุม $\hat{A}$ เป็นมุมฉาก ให้ $D$ เป็นจุดบนด้าน $BC$ ซึ่งทำให้ $AD$ ตั้งฉากกับ $BC$ ถ้าด้าน $AB$ ยาว 1 หน่วย และด้าน $AC$ ยาว 3 หน่วย จงหาค่าของ $\frac{CD}{BD}$
\end{minipage}
\hfill
\begin{minipage}{0.38\textwidth}
    \centering
    % บล็อกจำลองรูปภาพขนานไปกับข้อความด้านซ้าย
    \framebox[4cm]{\vbox to 3cm{\vfill\hbox to 4cm{\hfill [รูปด้านขวา] \hfill}\vfill}}
\end{minipage}
\par\medskip
\choiceshorizontal{3}{4}{9}{10}

\pagebreak
% ========= part 3 ============
\examsection{ตอนที่ 3 วิทยาการคำนวณ แบบอัตนัย (เติมคำตอบ) จำนวน 5 ข้อ (ข้อ 56--60) ข้อละ 2 คะแนน}

% ตัวอย่างโจทย์อัตนัย / เติมคำตอบ (ไม่ต้องมีช้อยส์)
\question บริษัทของคุณพ่อรับจัดสวนและปูสนามหญ้าที่บ้านแห่งหนึ่ง สนามหญ้ามีขนาด กว้าง 2 เมตร ยาว 20 เมตร คุณพ่อได้จัดวางสปริงเกอร์ไว้ 8 ตัว ซึ่งสปริงเกอร์แต่ละตัวมีขนาดไม่เท่ากัน สามารถฉีดน้ำได้รอบทิศตามรัศมีของแต่ละตัว โดยวางตรงเส้นประกึ่งกลางตามความยาวของสนามหญ้า จะต้องเปิดสปริงเกอร์อย่างน้อยที่สุดกี่ตัว จึงจะครอบคลุมพื้นที่สนามหญ้าทั้งหมด

\question หากกำหนดไวยากรณ์ (Grammar) ของภาษาหนึ่งเป็น $(01)^*00(11)^*$ ผลการรับรู้ (Recognition) ของสตริงคำว่า \texttt{010100} จะเสร็จสิ้นที่สถานะใด และสอดคล้องกับโครงสร้างไวยากรณ์หรือไม่

\end{document}